% TOP_PROBLEM: {"id": 1861, "title": "Erdős $5000 Problem on Prime Gaps", "category_id": 1, "category": {"id": 1, "name": "number_theory", "display_name": "Number Theory", "description": "Properties of integers, prime numbers, Diophantine equations.", "slug": "number-theory", "order_index": 1, "created_at": "2026-07-31T15:26:25.670Z"}, "status": "open"}
% FIRST_POSTED: null
% ENTRY_KIND: statement_only
% Imported from: batch08.tex; UnsolvedMath ID 1861
\documentclass[11pt]{article}
\usepackage[margin=25mm]{geometry}
\usepackage{fontspec}
\setmainfont{Latin Modern Roman}
\usepackage{amsmath,amssymb,mathtools}
\usepackage{unicode-math}
\setmathfont{Latin Modern Math}
\usepackage{xurl,hyperref}
\hypersetup{colorlinks=true,urlcolor=blue}
\setcounter{secnumdepth}{0}
\setlength{\parindent}{0pt}
\setlength{\parskip}{5pt}
\setlength{\emergencystretch}{3em}
\newcommand{\sref}[2]{\hyperlink{source:#1:#2}{[S#2]}}
\newcommand{\cataloguescope}[1]{\paragraph{Recorded target.}#1}
\begin{document}
% UnsolvedMath ID 1861; source catalogue code GUY-A8a
\setcounter{section}{1860}
\section{Erdős \$5000 Problem on Prime Gaps}\label{1861}
{\small\sffamily UnsolvedMath ID 1861 · Number Theory}
\subsection{Definitions and mathematical statement}

\paragraph{Shared notation.}$\mathbb N=\{1,2,\ldots\}$, and $\mathbb Z,\mathbb Q,\mathbb R,\mathbb C$ denote the integers, rationals, reals and complex numbers. Any other conventions are specified locally.


Let \(p_n\) be the \(n\)-th prime and
\(d_n=p_{n+1}-p_n\). Let \(\log\) denote the natural logarithm,
and define its iterates \(\log_1x=\log x\) and
\(\log_{j+1}x=\log(\log_jx)\), whenever these expressions exist.
For sufficiently large \(x\), put
\[
                     R(x)=\frac{\log x\,\log_2x\,\log_4x}
                                       {(\log_3x)^2}>0.
\]
This is the Rankin scale in the recorded question.
Ford, Green, Konyagin and Tao proved the arbitrary-constant growth
result in the cited theorem, settling this historical target.
\paragraph{Statement recorded in the supplied source.}
Is it true that, for every real constant \(c>0\) and every
integer \(N\ge1\), there is an integer \(n\ge N\), large enough
that \(R(n)\) is defined and positive, satisfying
\[
                         d_n>c\,R(n)?
\]
Equivalently, is \(\limsup_{n\to\infty}d_n/R(n)=\infty\)? \sref{1861}{2}
\subsection{Short English statement}
Do consecutive prime gaps exceed every fixed multiple of the displayed four-logarithm Rankin scale infinitely often? This historical question was answered affirmatively.
\subsection{Sources}
\begin{description}
\item[\hypertarget{source:1861:1}{S1}] ULAM.AI and the UnsolvedMath Contributors, \emph{UnsolvedMath: A Curated Collection of Open Mathematics Problems}, record 1861 (GUY-A8a). CC BY 4.0. \url{https://huggingface.co/datasets/ulamai/UnsolvedMath}.
\item[\hypertarget{source:1861:2}{S2}] Kevin Ford, Ben Green, Sergei Konyagin and Terence Tao, \emph{Large gaps between consecutive prime numbers}, Annals of Mathematics 183 (2016), 935--974; arXiv:1408.4505v2, Section 1, Theorem 1 and the arbitrary-constant Rankin-scale result. \url{https://arxiv.org/abs/1408.4505}.
\end{description}
\label{end:1861}
\end{document}
